% Supplementary elementary proof, using the notation and theorem labels of
% five-colour-cubes.tex. Not included in the main manuscript: the full
% classification now supplies this necessity implication there.
\begin{proof}[Elementary necessity for the three-variety classification]
\emph{Necessity.} Theorem~\ref{thm:multiplicity-five} forces multiplicities $(4,5,5)$;
write $A$ for the four-copy variety. Every target has eight positions,
whereas any two source varieties supply at least nine copies. A Hall-deficient
set can therefore have neighbours in at most one source variety. A target
is solvable precisely when every position has a supplier and at most
$4,5,5$ positions are supplied exclusively by $A,B,C$, respectively.
In particular, failure to build the three sources themselves implies
\begin{align*}
 |K(A)\cap(K(B)\cup K(C))|&\le3,\\
 |K(B)\cap(K(A)\cup K(C))|&\le2,\\
 |K(C)\cap(K(A)\cup K(B))|&\le2.
\end{align*}
No pair is disjoint, by Lemma~\ref{lem:disjoint-pair}, so each pairwise
intersection has size one or two.

\emph{Possible source types.} No two adjacent varieties have the same sign. Otherwise, after reversal,
their cuts are $C(uv),C(uw)$, with intersection $\{ux,uy\}$ on the remaining
vertices $x,y$. At least one, say $X$ with cut $C(uv)$, has five copies. Its intersection
with the third variety $Z$ must be a nonempty subset of $\{ux,uy\}$.
If $Z$ is positive, its cut must be $C(uw)$, giving six shared corners with
the second variety. If $Z$ has negative cut $C(S)$ and contains, say, $ux$, the bound
two excludes $v$ from $S$; since $x\notin S$, this forces $S=\{u,w,y\}$,
which shares four arcs with $C(uw)$. If $Z$ is opposite, it contains both
$ux,uy$ and has parts $\{x,y\},\{u,z\}$, with $z=v$ or $w$; again it
shares four arcs with one of the two cuts. All are contradictions.

Three adjacent varieties are excluded by this observation. Two opposite
varieties and one adjacent are also impossible: their pairwise intersections
would coincide in a bidirected edge, which a cut cannot contain. If all
three are opposite, their undirected four-cycles share one edge and are
otherwise pairwise edge-disjoint, so their union is $K_5$. Choose an opposite
target avoiding one endpoint of the common edge. Each required edge then
belongs to exactly one source, and each source supplies at most two target
edges, hence at most four directed corners. Its copies suffice.

\emph{Normal form.} Thus exactly one variety is opposite. It cannot have five copies, for its
two overlaps would coincide in two tricoloured corners, forcing the adjacent
cuts to have the same sign. Consequently $A$ is opposite and $B,C$ have
opposite signs. Interchange $B,C$ if necessary so that $B$ is positive.
Their intersections with $A$ have size two, while
$K(B)\cap K(C)\subseteq K(A)$ is nonempty and, by the table, a single arc.
Write the smaller pairs as $\{o,u\},\{o,v\}$; the common arc is $uv$.
The four-cycle of $A$ must omit $o$: otherwise one of these pairs would
be a whole part of its $K_{2,2}$, giving four shared arcs. Normalize
$o=0,u=1,v=3$ and name their partners $2,4$. This gives the stated cuts,
with $r\in\{2,3,4\}$, $s\in\{1,2,4\}$ and $r\ne s$; equality would
add the common bicoloured corner $b(r,0)$.

\end{proof}
